\section{Background and Related Work}

Causal reasoning goes beyond associational prediction by asking how a system would change under interventions and counterfactual alternatives \citep{Pearl_2009,pearl2018why,imbens_rubin_causal_2015}. Structural causal models (SCMs) formalize these assumptions as directed graphs plus structural equations \citep{Pearl_2009}. In \projectname{}, each episode's hidden SCM is both the ground truth (\S\ref{sec:scm-construction}) and the evaluation target (\S\ref{sec:eval-targets}), letting us score whether an agent recovers the graph and target equation, not only whether it predicts the reactor value.

Most LLM causal evaluations are static: they ask models to answer textual causal questions, reason over given graphs, classify cause--effect direction, or solve formal causal-inference queries \citep{kiciman2023causal,jin2024cladderassessingcausalreasoning,corr2cause2024,clear2024,wang-2024-causalbench,chen2024causal}. Related work also uses LLMs as causal priors for edge scoring, causal ordering, or query-efficient discovery \citep{llmcausalgraphs2023,llmeffectivepriors2024,causalorder2023,efficientcausalgraph2024}. Recent SCM-oriented studies either use LLM metadata reasoning to support graph discovery \citep{abdulaal2024causalmodellingagents} or test coefficient elicitation when the DAG is supplied \citep{yamaoka2026linearllmscm}. HypoBench further shows that hypothesis-generation benchmarks must account for how prior knowledge shapes model behavior \citep{liu2025hypobench}. These settings clarify what causal knowledge LLMs can express, but they usually provide the variables, evidence, graph, or query up front. \projectname{} instead asks whether an LLM agent can gather evidence, revise a hypothesis, and transfer the learned mechanism to a new instance, all within a scientific-discovery setting that offers no hints about the underlying causal structure.

Interactive environments broaden evaluation beyond one-shot answers, including scientific-discovery worlds, budgeted graph-discovery games, causal games, and non-LLM intervention planners \citep{DBLP:journals/corr/abs-2406-06769,igda2025,chen2026causalgame,gregorini2025dodo}. A basic agent scaffold for such settings is ReAct-style reasoning and acting, where the model interleaves deliberation with executable environment actions \citep{DBLP:conf/iclr/YaoZYDSN023}. The closest recent benchmark is Auto-Bench, where LLM agents iteratively query scientific or social-network environments to recover a hidden causal graph \citep{chen2025autobench}. Work on black-box reverse engineering similarly shows that actively designing queries is not equivalent to receiving another agent's intervention data \citep{geng2025reliableaiscientists}.

\projectname{} differs from Auto-Bench in its evaluation target. Auto-Bench primarily asks whether an agent can discover a hidden DAG through interaction. \projectname{} asks whether the discovered mechanism \emph{transfers}: after learning from prior measurements and interventions on a manipulator crystal, the agent must predict a held-out reactor crystal generated by the same SCM, while its per-step DSL hypotheses expose the graph, the \texttt{frequency} structural equation, and the coefficients it is committing to. This makes it possible to separate task utility from structural and quantitative faithfulness, and to audit how an LLM agent revises or fails to revise an explicit SCM over time.

This connects two evaluation practices: explicit SCM recovery from causal discovery and sequential experiment design from agent benchmarks. Because each episode has a known ground-truth SCM and a logged interaction trace, \projectname{} can score both final-task utility and the faithfulness of the recovered mechanism.

